% !PRINTED 334
\weberpage{334}
Denn die notwendige und hinreichende Bedingung, daß (2) Basis von $\mathfrak{a}$ sei, besteht darin, daß jede durch $\mathfrak{a}$ teilbare ganze Zahl $\alpha$ und nur solche in der Form:
\weq{3}{\alpha = x_1\alpha_1 + x_2\alpha_2 + \cdots + x_n\alpha_n}
mit ganzzahligen $x_1, x_2, \ldots, x_n$ enthalten ist. Folglich sind alle Zahlen
\weq{4}{\beta = x_1\alpha_1\eta + x_2\alpha_2\eta + \cdots + x_n\alpha_n\eta}
durch $\mathfrak{b}$ teilbar. Ist umgekehrt $\beta$ eine durch $\mathfrak{b}$ teilbare ganze Zahl, so ist $\beta/\eta$ durch $\mathfrak{a}$ teilbar und folglich in der Form (3) darstellbar. Mithin ist $\beta$ durch (4) darstellbar.

Ferner ist nach der Bezeichnung von (9) und (10), § 87, wenn $\alpha$ eine positive Basis ist:
\[
\begin{aligned}
\Sigma \pm \alpha_{1,1}\alpha_{2,2}\ldots\alpha_{n,n} &= N(\mathfrak{a})\, \Sigma \pm \omega_{1,1}\omega_{2,2}\ldots\omega_{n,n}, \\
\Sigma \pm \beta_{1,1}\beta_{2,2}\ldots\beta_{n,n} &= N(\eta)\, \Sigma \pm \alpha_{1,1}\alpha_{2,2}\ldots\alpha_{n,n}, \\
&= N(\eta)N(\mathfrak{a})\, \Sigma \pm \omega_{1,1}\omega_{2,2}\ldots\omega_{n,n}, \\
&= \pm N(\mathfrak{b})\, \Sigma \pm \omega_{1,1}\omega_{2,2}\ldots\omega_{n,n},
\end{aligned}
\]
und da nach (1)
\weq{5}{N(\mathfrak{b}) = N(\eta)N(\mathfrak{a})}
ist, so muß bei $\pm N(\mathfrak{b})$ das positive Zeichen stehen.

Ist $\lambda$ eine Basisform von $\mathfrak{a}$, so ist $\eta\lambda$ eine Basisform von $\mathfrak{b}$, und
\weq{6}{\begin{aligned} N(\lambda) &= N(\mathfrak{a})T, \\ N(\eta\lambda) &= N(\mathfrak{b})T, \end{aligned}}
und zu den beiden Idealen $\mathfrak{a}$, $\mathfrak{b}$ gehört also dieselbe Form $T$. Vereinigen wir äquivalente Formen $T$ in eine Klasse, so können wir nach 2. sagen:

3. \gesperrt{Zu jeder Idealklasse gehört eine bestimmte Formenklasse.}

Bei der Beantwortung der Frage, ob zu einer und derselben Formenklasse verschiedene Idealklassen gehören können, machen wir die Voraussetzung, daß $\Omega$ ein \gesperrt{Normalkörper} sei: Nehmen wir also an, zu zwei Idealen $\mathfrak{a}$ und $\mathfrak{a}'$ gehören äquivalente Formen $T$ und $T'$, so ist, wenn $\lambda$ eine Basisform von $\mathfrak{a}$ ist,
\weq{7}{N(\lambda) = N(\mathfrak{a})T,}
und wenn wir $T'$ durch eine lineare Substitution in $T$ überführen, und unter $\lambda'$ eine Basisform von $\mathfrak{a}'$ (mit denselben Variablen $t$ wie $\lambda$) verstehen, so ist
\weq{8}{N(\lambda') = N(\mathfrak{a}')T.}
